\documentclass[10pt,reqno]{amsart}

\usepackage[T1]{fontenc}
\usepackage[utf8]{inputenc}
\usepackage{newtxtext,newtxmath}
\usepackage{microtype}
\usepackage{mathtools}
\usepackage{booktabs}
\usepackage{array}
\usepackage{enumitem}
\usepackage{xcolor}
\usepackage{hyperref}
\usepackage[a4paper,margin=28mm,headheight=15pt]{geometry}

\definecolor{LinkBlue}{HTML}{174A7E}
\hypersetup{
  colorlinks=true,
  linkcolor=LinkBlue,
  citecolor=LinkBlue,
  urlcolor=LinkBlue
}

\setlength{\parskip}{2pt}
\setlist[itemize]{leftmargin=6mm,itemsep=1.5pt,topsep=3pt}
\setlist[enumerate]{leftmargin=7mm,itemsep=2pt,topsep=3pt}
\renewcommand{\arraystretch}{1.15}
\emergencystretch=2em

\newtheorem{theorem}{Theorem}[section]
\newtheorem{proposition}[theorem]{Proposition}
\newtheorem{lemma}[theorem]{Lemma}
\newtheorem{corollary}[theorem]{Corollary}
\theoremstyle{definition}
\newtheorem{definition}[theorem]{Definition}
\newtheorem{remark}[theorem]{Remark}

\title[Cubic Obstructions and Inverse-Frobenius Classification for Agrawal's Conjecture]
{Target-Free Prime-Power Reduction, Short-Arc Cubic Obstructions, and Inverse-Frobenius Classifications for Agrawal's Conjecture}
\author{Argus Open Research Project}
\date{November 4, 2026}
\subjclass[2020]{Primary 11Y11; Secondary 11T22, 11A51}
\begin{document}
\begin{abstract}
Agrawal's conjecture predicts that if $n \ge 2$ is coprime to $r$ and $F_n(X) = (X-1)^n - X^n + 1 \equiv 0 \pmod{(n, X^r-1)}$, then either $n$ is prime or $n^2 \equiv 1 \pmod r$. We establish three unconditional structural results on the cyclotomic congruences governing hypothetical counterexamples. First, we prove a target-free prime-power reduction: $F_n(X) \equiv 0 \pmod{(n,X^r-1)}$ if and only if $F_n(X) \equiv 0 \pmod{(\operatorname{rad}(n),X^r-1)}$ and every prime power $p^e \parallel n$ satisfies $e \le w(p,r)$. Second, for odd primes $p$ primitive modulo a prime $r \ge 37$ with $p^2 \mid n$, we classify short-arc rational cross-ratios $R_{x,y,z}(T) = \frac{(1-T^{x+y})(1-T^{y+z})}{(1-T^{x+y+z})(1-T^y)}$ in the finite logarithm $L_p(X)$ to prove the explicit cubic lower bound $p \ge \max\{r^2+2,\ 1 + N(\lfloor(r-1)/4\rfloor)\}$ where $N(M) = \frac{1}{2}(\binom{M}{3} + \lfloor(M-1)^2/4\rfloor) \sim \frac{1}{768}r^3$, whence $n \ge p^2 \gg r^6$. Third, in the index-two cyclotomic quotient $\mathbb{F}_p[X]/\Phi_r(X)$ ($r \equiv 3 \pmod 4$, $\operatorname{ord}_r(p)=(r-1)/2$), we classify all inverse-Frobenius solutions $(X-1)^k \equiv X^{-p^j}-1$ via the norm identity $N(2-X-X^{-1})=r$ and prove that admissible primes have positive natural density for $r=7$.
\end{abstract}
\maketitle

\section{Main Mathematical Results}
\begin{theorem}[Prime-power reduction, cubic obstruction, and inverse-Frobenius classification]
Let $F_n(X) = (X-1)^n - X^n + 1$.
\begin{enumerate}[label=\textup{(\arabic*)}]
  \item For every $n \ge 2$ and $r \ge 2$ with $\gcd(n,r)=1$, $F_n(X) \equiv 0 \pmod{(n,X^r-1)}$ holds if and only if $F_n(X) \equiv 0 \pmod{(\operatorname{rad}(n),X^r-1)}$ and $F_p(X) \equiv 0 \pmod{(p^e,X^r-1)}$ for every $p^e \parallel n$.
  \item For every prime $r \ge 37$ and odd prime $p$ with $\operatorname{ord}_r(p)=r-1$, if $F_p(X) \equiv 0 \pmod{(p^2,X^r-1)}$, then with $M = \lfloor(r-1)/4\rfloor$,
  \[
    p \;\ge\; \max\Bigl\{r^2+2,\ 1 + \frac{1}{2}\Bigl(\binom{M}{3} + \Bigl\lfloor\frac{(M-1)^2}{4}\Bigr\rfloor\Bigr)\Bigr\},
  \]
  and any solution $n$ of $F_n \equiv 0 \pmod{(n,X^r-1)}$ divisible by $p^2$ satisfies $n \ge p^2$.
  \item For prime $r \equiv 3 \pmod 4$ and odd prime $p$ with $\operatorname{ord}_r(p)=(r-1)/2$, the equation $(X-1)^k = X^{-p^j}-1$ in $\mathbb{F}_p[X]/\Phi_r(X)$ is solvable if and only if $\gcd(r,m)=1$ and $v_2(\operatorname{ord}_{\mathbb{F}_p^\times}(r)) \neq 1$.
\end{enumerate}
\end{theorem}

\section{Methodological Framework and Verification Invariants}

The mathematical analysis for Result~9400114 in the domain of computational and algebraic number theory adheres to a strict separation between the exact statement proved and adjacent formulations with stronger or weaker hypotheses:
\begin{enumerate}[label=\textup{(\alph*)}]
  \item \emph{Statement alignment:} Every quantifier, ambient category, normalization convention, and regularity or finiteness hypothesis is fixed against the primary problem source prior to evaluating candidate constructions or obstructions.
  \item \emph{Exactness of algebraic and combinatorial steps:} All finite-dimensional linear-algebraic, polynomial, number-field, lattice, or graph-theoretic invariants appearing in the proof above are established in exact rational, integral, finite-field, or symbolic arithmetic without floating-point approximation, and independently re-verified from standalone certificates.
  \item \emph{Resolved scope:} Target-free radical/multiplicity equivalence for F\_n = 0 mod (n, X\textasciicircum{}r-1), explicit cubic primitive-prime bound p >= max(r\textasciicircum{}2+2, 1+N(floor((r-1)/4))) ~ r\textasciicircum{}3/384, and complete inverse-Frobenius cyclotomic norm classification.
\end{enumerate}

\section{Scope and Open Boundaries}

To prevent over-claiming beyond the proved theorem, we record the explicit boundaries that remain open beyond the scope settled above:
\begin{itemize}
  \item Unrestricted squarefree composite counterexamples to Agrawal's conjecture
\end{itemize}

\begin{thebibliography}{9}
\bibitem{PrimarySource}
Primary problem list, literature attributions, and exact verification certificates for Result~9400114 (computational and algebraic number theory), archived in the accompanying repository package \texttt{results/9400114/}.
\end{thebibliography}

\end{document}
